%--------------------------------------------------------------------------------------
% Syntax highlighting
%--------------------------------------------------------------------------------------
% Two packages, one look:
%  - minted (Pygments) for common languages: \begin{minted}{python} ... \end{minted}
%    inside a listing float, or \mintinline{rust}{...} in text. Needs shell escape,
%    which Overleaf enables; locally compile with -shell-escape.
%  - listings for languages Pygments does not know: Refinery, Tamarin, Lean.
%    \begin{lstlisting}[style=refinery, caption=..., label=lst:...] ... \end{lstlisting}
% Both are numbered "Listing N" in one sequence and both work with \cref.

% Colors of the Pygments default style, so listings and minted code look alike.
\definecolor{codeKeyword}{HTML}{008000}
\definecolor{codeType}{HTML}{B00040}
\definecolor{codeString}{HTML}{BA2121}
\definecolor{codeComment}{HTML}{3D7B7B}
\definecolor{codeEmph}{HTML}{0000FF}

% minted: the same frame, background and size as the listings in include/preamble.tex.
\setminted{
  style=default,
  fontsize=\scriptsize,
  frame=single,
  bgcolor=lightgray,
  breaklines=true,
  tabsize=2,
}
\setmintedinline{bgcolor={}}

% One counter for lstlisting and minted's listing float.
\makeatletter
\AtBeginDocument{%
  \let\c@listing\c@lstlisting
  \renewcommand{\thelisting}{\thelstlisting}%
}
\makeatother
\crefname{listing}{Listing}{Listings}

%--------------------------------------------------------------------------------------
% Refinery (https://refinery.tools), after the SoSyM 2026 paper's definition.
% Add the metamodel's own type names per listing: [style=refinery, morekeywords=[2]{Credential}]
\lstdefinelanguage{refinery}{
  morekeywords={class,abstract,extends,contains,refers,opposite,container,
    enum,pred,error,shadow,propagation,rule,pattern,
    must,may,true,false,unknown,count,exists,equals,default,
    scope,node,import,int,real,string,boolean},
  morekeywords=[2]{},
  alsoletter={_},
  sensitive=true,
  morecomment=[l]{//}, morecomment=[s]{/*}{*/}, morecomment=[l]{\%},
  morestring=[b]{"},
}
\lstdefinestyle{refinery}{
  language=refinery,
  keywordstyle=\color{codeKeyword}\bfseries,
  keywordstyle=[2]\color{codeType}\bfseries,
  literate=%
    {==>}{$\Longrightarrow$}3
    {<->}{$\longleftrightarrow$}3
    {<=}{$\leqslant$}1
    {>=}{$\geqslant$}1
    {!=}{{$\neq$}}1,
}

%--------------------------------------------------------------------------------------
% Tamarin prover theories (.spthy), after Smuk Ferenc's TDK definition.
\lstdefinelanguage{tamarin}{
  morekeywords={theory,begin,end,builtins,functions,equations,rule,lemma,restriction,
    let,in,All,Ex,not,exists-trace,all-traces,reuse,use_induction,sources,
    heuristic,tactic,predicates,macros,private},
  morekeywords=[2]{Fr,In,Out,K},
  alsoletter={-},
  sensitive=true,
  morecomment=[l]{//}, morecomment=[s]{/*}{*/},
  morestring=[b]{"},
}
\lstdefinestyle{tamarin}{
  language=tamarin,
  keywordstyle=\color{codeKeyword}\bfseries,
  keywordstyle=[2]\color{codeType}\itshape,
  literate=%
    {-->}{$\longrightarrow$}3
    {==>}{$\Longrightarrow$}3
    {--[}{{-\hspace{-0.2em}-[}}3
    {]->}{{]\hspace{-0.2em}$\rightarrow$}}3
    {@}{{$@$}}1
    {<=>}{$\Leftrightarrow$}3,
}

%--------------------------------------------------------------------------------------
% Lean 4. Unicode notation is typeset as math, so it renders with every engine.
% Only in block listings: \lstinline with Unicode characters fails under XeLaTeX.
\lstdefinelanguage{lean}{
  morekeywords={theorem,lemma,def,example,abbrev,structure,inductive,class,instance,
    where,deriving,namespace,section,end,open,variable,import,universe,noncomputable,
    fun,match,with,if,then,else,let,have,show,from,by,at,calc,do,return,
    Type,Prop,Sort},
  morekeywords=[2]{intro,intros,exact,apply,refine,rw,simp,simp_all,omega,decide,
    cases,induction,rcases,obtain,constructor,unfold,contradiction,assumption,
    exfalso,left,right,use,linarith,norm_num,aesop,trivial,rfl,sorry},
  alsoletter={_'.},
  sensitive=true,
  morecomment=[l]{--}, morecomment=[n]{/-}{-/},
  morestring=[b]{"},
}
% XeLaTeX and LuaLaTeX print Lean's Unicode directly in DejaVu Sans Mono, which has the
% glyphs (Latin Modern Mono does not); pdfLaTeX typesets them as math instead.
\ifxetexorluatex
\newfontfamily\leanfont{DejaVuSansMono.ttf}[Scale=MatchLowercase, BoldFont=DejaVuSansMono-Bold.ttf,
  ItalicFont=DejaVuSansMono-Oblique.ttf, BoldItalicFont=DejaVuSansMono-BoldOblique.ttf]
\lstdefinestyle{lean}{
  language=lean,
  basicstyle=\scriptsize\leanfont,
  keywordstyle=\color{codeKeyword}\bfseries,
  keywordstyle=[2]\color{codeEmph},
  literate={},
}
\else
\lstdefinestyle{lean}{
  language=lean,
  keywordstyle=\color{codeKeyword}\bfseries,
  keywordstyle=[2]\color{codeEmph},
  literate=%
    {∀}{{$\forall$}}1 {∃}{{$\exists$}}1 {λ}{{$\lambda$}}1 {¬}{{$\neg$}}1
    {→}{{$\to$}}1 {←}{{$\leftarrow$}}1 {↔}{{$\leftrightarrow$}}1
    {∧}{{$\land$}}1 {∨}{{$\lor$}}1 {≤}{{$\le$}}1 {≥}{{$\ge$}}1 {≠}{{$\ne$}}1
    {∈}{{$\in$}}1 {∉}{{$\notin$}}1 {⊆}{{$\subseteq$}}1 {∪}{{$\cup$}}1 {∩}{{$\cap$}}1
    {⟨}{{$\langle$}}1 {⟩}{{$\rangle$}}1 {×}{{$\times$}}1 {·}{{$\cdot$}}1
    {ℕ}{{$\mathbb{N}$}}1 {ℤ}{{$\mathbb{Z}$}}1 {⊢}{{$\vdash$}}1 {▸}{{$\triangleright$}}1,
}
\fi
